\documentclass{article}
\usepackage[utf8]{inputenc}
\usepackage{amsmath,amssymb}
\usepackage[most]{tcolorbox}
\usepackage{geometry}
\geometry{margin=2.5cm}

\newcommand{\step}[1]{\par\noindent\hangindent=3.5em\hangafter=1%
  \makebox[3.5em][l]{\textbf{#1.}}}
\newcommand{\steptag}[1]{\hfill{\small\textsf{[#1]}}}

\newtcolorbox{proofbox}[1]{
  colback=gray!5, colframe=gray!60,
  fonttitle=\bfseries, title={#1},
  breakable, enhanced,
  left=4pt, right=4pt, top=4pt, bottom=4pt,
  before skip=10pt, after skip=10pt
}

\begin{document}

\begin{proofbox}{Cohomological Kunneth isomorphism over R: for CW complexe...}
\step{1} Cohomological Kunneth isomorphism over R: for CW complexes X and Y with each H\textasciicircum{}k(Y;R) a finitely generated free R-module, the cross product H*(X;R) (x)\_R H*(Y;R) -> H*(X x Y;R) is an isomorphism of rings; in particular this holds over the field R = reals for finite-CW spaces with finite-dimensional cohomology. \steptag{validated} \\[6pt]
\step{1.1} By the allowed external GT-hatcher-AT-thm-3.15 (Hatcher, Theorem 3.15, refs/hatcher-algebraic-topology/AT.txt:13505-13506), for an arbitrary commutative coefficient ring R, if X and Y are CW complexes and H\textasciicircum{}k(Y;R) is a finitely generated free R-module for every k, then the cross product H*(X;R) (x)\_R H*(Y;R) -> H*(X x Y;R) is an isomorphism of rings. \steptag{validated} \\[6pt]
\step{1.1.1} Scope correction and application: in node 1.1 and in the general clause of root node 1, R ranges only over commutative coefficient rings, exactly the ambient class introduced in the source immediately before GT-hatcher-AT-thm-3.15; no assertion is made for noncommutative rings. With this restriction, the hypotheses that X and Y are CW complexes and every H\textasciicircum{}k(Y;R) is finitely generated free are precisely those of GT-hatcher-AT-thm-3.15, whose conclusion is the displayed cross-product ring isomorphism. The later specialization remains valid because the real numbers form a commutative field. \steptag{validated} \\[4pt]
\step{1.2} In the stated special case R = real numbers, with X and Y finite-CW spaces and H\textasciicircum{}k(Y;R) finite-dimensional for every k, X and Y are CW complexes and each H\textasciicircum{}k(Y;R) is a finitely generated free R-module. \steptag{validated} \\[6pt]
\step{1.2.1} By definition, a finite-CW space is a CW complex having finitely many cells; hence the assumed finite-CW spaces X and Y satisfy the CW-complex hypothesis of GT-hatcher-AT-thm-3.15. \steptag{validated} \\[4pt]
\step{1.2.2} For every k, finite-dimensionality of the real vector space H\textasciicircum{}k(Y;real numbers) means it has a finite basis. The same basis makes H\textasciicircum{}k(Y;real numbers), viewed as a module over the field of real numbers, free and generated by finitely many elements. Thus it is a finitely generated free real module for every k. \steptag{validated} \\[4pt]
\step{1.3} Applying the general isomorphism in node 1.1 using the verified hypotheses in node 1.2 yields the particular assertion: over R = real numbers for finite-CW spaces with finite-dimensional cohomology, the cross product is an isomorphism of rings. \steptag{validated} \\[4pt]
\end{proofbox}

\end{document}
